\documentclass[11pt]{article}

\usepackage[a4paper,margin=2.7cm]{geometry}
\usepackage{amsmath,amssymb,mathtools}
\usepackage{booktabs,longtable,array}
\usepackage{microtype}
\usepackage{enumitem}
\usepackage{xcolor}
\usepackage{xurl}
\usepackage[hidelinks]{hyperref}
\usepackage[T1]{fontenc}
\usepackage{lmodern}

\setlength{\parindent}{0pt}
\setlength{\parskip}{0.55em}
\setlist{nosep,leftmargin=*}
\providecommand{\tightlist}{\setlength{\itemsep}{0pt}\setlength{\parskip}{0pt}}

\title{\textbf{From Self-Maintenance to Evolvability:\\Organizational Accessibility and a Dynamical Ratchet for Cumulative Organization}}
\author{}
\date{Revised working paper --- 6 September 2026}

\begin{document}
\maketitle

\begin{abstract}
Why can organization accumulate rather than repeatedly being
rediscovered from scratch? We propose that the relevant object is not
organization alone, but the \textbf{accessibility of future
organization}: the probability that persistent successor organizations
can be reached from what currently exists within a finite horizon.
Realized organization can itself become a causal determinant of this
accessibility --- not merely influence the next state, which any
dynamical system does, but alter the transition structure through which
viable successor states become reachable.

We construct a minimal account in four steps. First, an open production
network can collectively replace the processes that constitute it,
yielding an organization reproduction number, \(R_{\rm org}\), formally
identical to a next-generation-matrix threshold: below one, organization
introduced at low abundance decays; above one, it can establish. Second,
a seed-dependent feedback can make an expanded organization easier to
maintain than to establish. In the worked model, a rare process \(C\)
cannot invade the resident \(A\)--\(B\) organization below
\(J_{\rm inv}=1.0667\), yet once established the expanded
\(A\)--\(B\)--\(C\) organization persists down to
\(J_{\rm fold}=0.5978\). This produces bistability and hysteresis: at
the same external environment, accessibility depends on the realized
state and its basin of attraction. Third, in reproducing systems,
retained organization becomes the inherited starting point of subsequent
variation. Descendants explore modifications of an already viable
structure rather than reconstructing viability de novo. Near a
persistence boundary such as \(R_{\rm org}=1\), comparatively small
adaptive changes may restore positive growth while the declining lineage
still generates a finite number of reproductive and exploratory trials.
This connects cumulative organization to inheritance, evolutionary
rescue, and eco-evolutionary feedback. Fourth, architectures may
themselves alter the distribution from which future variants are
generated. If such variation-generating mechanisms are heritable and
remain coupled to persistence or reproduction, the search process itself
can evolve --- an evolvability feedback, which existing theory shows to
be possible but conditional, not automatic.

We formalize the common object as a finite-horizon accessibility kernel
\(\mathcal A_T(S\mid\theta,z,E,q)\): the probability that a persistent
successor organization in a set \(S\) is realized within time \(T\),
conditional on current organization \(\theta\), its realized dynamical
state \(z\), external environment \(E\), and the variant-generating
mechanism \(q\). The mechanisms are mathematically distinct but share
one structure: realized organization changes this kernel. It can change
\textbf{what can remain}, \textbf{where future search begins and how
many exploratory trials occur}, and \textbf{how candidate successors are
generated}. This does not imply a force toward complexity or
monotonically increasing organization. Accessibility may narrow as well
as expand, and retained structures may be maladaptive. The framework
instead provides a non-teleological route by which previous viability
can become historically consequential: previous solutions can become
part of the conditions from which subsequent solutions are generated.
The central claim is therefore not that existence makes every future
existence easier, but that \textbf{existence changes what existence can
follow}; cumulative organization is the special case in which enough
previous viability is retained that subsequent viable states need not
repeatedly be reconstructed from scratch.
\end{abstract}

\section{The problem: not survival of parts, but continuation of
organization}
Organized systems persist despite continuous turnover. Molecules decay,
cells die, organisms replace cells, institutions replace members,
machines replace components, and computational processes terminate and
restart. Persistence therefore need not mean persistence of the same
material parts. It can instead mean persistence of a causal
organization: a network of processes whose activity repeatedly produces
the processes required for the network to continue.

This shifts the explanatory target. The question is not merely why a
component survives, but how a collection of locally acting processes can
collectively cause its own continuation. Everything that follows in this
paper is an attempt to answer one more general question that sits behind
that one, which we make the paper's central object rather than a
conclusion reached late: not how the present state influences the next
state --- every dynamical system does that --- but

\begin{quote}
\textbf{how realized organization can endogenously change the
accessibility of its own successors.}
\end{quote}

The argument proceeds in four steps, corresponding to a reproduction
threshold plus three distinct mechanisms by which realized organization
changes what is accessible next. The first step is addressed by a
reproduction threshold for organization, which sets up the ``remain or
not'' question that the three mechanisms subsequently modify. The first
mechanism, retention, is addressed by a seed-dependent nonlinear
extension that produces bistability and hysteresis (Section 6). The
second, inheritance, is addressed by a rescue-type argument transplanted
from population biology into the same reproduction-number formalism
(Section 7). The third, evolvability, is presented as an explicit
extension of the variant-generating mechanism itself (Section 8) --- a
hypothesis, not a result already proved by the first two.

The construction sits near several established traditions.
Autocatalytic-set theory studies networks in which reactions are
collectively sustained by catalysts produced within the network (Hordijk
et al., 2010, 2012). Jain and Krishna showed how autocatalytic structure
can emerge and reorganize in evolving catalytic networks (Jain and
Krishna, 1998, 2001). Peng, Linderoth and Baum formalized seed-dependent
autocatalytic systems that can remain active after a transient seed and
can form hierarchical structures (Peng et al., 2022).
Evolutionary-rescue theory studies how populations facing environmental
deterioration can avoid extinction through adaptation before decline is
complete (Bell, 2017; Ferriere and Legendre, 2013), and
pathogen-emergence theory has shown that a reproduction number near its
threshold, rather than far below it, is precisely where transmission
chains lengthen and the opportunity for adaptive mutation grows (Antia
et al., 2003; Woolhouse et al., 2005). The present contribution is
narrower: we use an explicit resource balance, derive a
reproduction-number threshold, show how one added feedback generates a
backward-bifurcation geometry, show that a second and independent
mechanism operates near that same threshold without requiring the
feedback loop at all, and then use both, plus a third hypothesized
mechanism, to define what accessibility means for a self-maintaining
organization.

\subsection{Organizational
accessibility}
The common object underlying the mechanisms developed below is
\textbf{organizational accessibility}. To keep architecture distinct
from instantaneous state, let \(\theta\) denote organizational structure
--- for example, which processes and interactions are available --- and
let \(z\) denote the current dynamical realization of that structure,
such as the abundance vector \((a,b,c,r)\). Let \(E\) denote the
external environment (here including \(J,\ell,d\)), and let \(q\) denote
the mechanism by which candidate successor organizations are generated.

Rather than using ``accessibility'\,' interchangeably for probability,
rate, cost, and difficulty, we define a finite-horizon transition
kernel. For a set \(S\) of candidate successor organizations,

\[
\boxed{
\mathcal A_T(S\mid\theta,z,E,q)
=
\Pr\!\left(\text{a persistent successor in }S\text{ is realized by time }T\mid\theta,z,E,q\right)
}
\tag{0}
\]

The kernel may of course be induced by underlying rates, energetic or
material costs, mutation or recombination probabilities, developmental
constraints, and basin geometry. Those quantities are mechanisms
determining \(\mathcal A_T\); they are not alternative meanings of it.

The central claim of this paper is not that organization necessarily
increases, nor that successful systems anticipate useful futures. It is
that \textbf{realized organization can itself become a causal
determinant of the transition kernel over future viable organization}.
The three mechanisms developed below modify that kernel through
different causal channels:

\begin{itemize}
\tightlist
\item
  \textbf{Mechanism I --- retention} (Section 6). At identical external
  conditions, different realized states can lie in different basins of
  attraction. An established organization can therefore remain
  accessible where the same organization could not be established from
  rarity. Realization changes \textbf{what can remain}.
\item
  \textbf{Mechanism II --- inherited adaptive extension} (Section 7).
  Reproduction with inheritance makes the current organization the base
  point of future variation rather than the organization-free state.
  Continued reproduction also supplies a finite number of exploratory
  trials before extinction. Realization changes \textbf{where the next
  search begins and how many attempts are available}.
\item
  \textbf{Mechanism III --- evolvability} (Section 8). Organization can
  alter \(q\), the mechanism that generates, connects, tests, or
  preserves candidate successors. Realization changes \textbf{how the
  next search is performed}.
\end{itemize}

None of these mechanisms need favour greater complexity, adaptability,
or persistence. They may preserve maladaptive states, constrain
exploration, remove functions, or favour simplification. Their common
feature is descriptive rather than normative:

\[
\boxed{\text{Realized organization transforms the accessibility of organization that can follow it.}}
\]

The transformation is recursive because a successor that persists
becomes the current organization and state from which later
accessibility is determined:

\[
(\theta_t,z_t)
\longrightarrow
\mathcal A_T(\cdot\mid\theta_t,z_t,E_t,q_t)
\longrightarrow
(\theta_{t+1},z_{t+1})
\longrightarrow
\mathcal A_T(\cdot\mid\theta_{t+1},z_{t+1},E_{t+1},q_{t+1})
\longrightarrow\cdots.
\]

History therefore changes not merely the state of the system, but the
transition structure available from that state. We define a
\textbf{cumulative organizational ratchet} as a history-dependent
transformation of this accessibility kernel in which enough previously
achieved viability is retained that subsequent viable organization need
not be rediscovered de novo.

Sections 2--3 derive the reproduction threshold determining whether
organization is realized at all. Sections 6--8 then develop three
distinct mechanisms by which realized organization feeds back into
future accessibility.

\section{A resource-limited self-maintaining production
network}
Let \(x_i \geq 0\) denote the abundance of executing process type \(i\),
and let \(r \geq 0\) denote activated substrate. The nonnegative matrix
\(B\) contains local production relations: \(B_{ij}\) is the rate at
which process \(j\) produces process \(i\) per unit substrate and
producer abundance. Each new unit consumes one unit of substrate.
Existing processes are lost at rate \(d > 0\). The environment supplies
activated substrate at rate \(J\), while unused substrate is lost at
rate \(\ell r\). The model is

\[
\dot{x} = rBx - dx, \tag{1}
\] \[
\dot{r} = J - \ell r - r\mathbf{1}^T Bx. \tag{2}
\]

The same production term enters positively in the process balance and
negatively in the substrate balance. Hence this is a replacement model
rather than a model in which producers are immortal. Defining

\[
T = r + \sum_i x_i,
\]

we obtain

\[
\dot{T} = J - \ell r - d\sum_i x_i \leq J - \min(\ell, d)T.
\]

Thus total activated material remains bounded, and when \(J = 0\) the
active organization ultimately disappears. The organization is causally
closed only in the limited sense that its processes can regenerate one
another; energetically and materially it remains open.

This distinction is important. The model is not a derivation from full
chemical thermodynamics. A real molecular implementation would require
stoichiometric, kinetic and free-energy constraints to be specified
separately. Thermodynamic work on molecular replicators makes clear that
replication and selection carry energetic requirements rather than
escaping them (Kolchinsky, 2025).

\section{A reproduction number for
organization}
With no active organization, the substrate converges to

\[
r_0 = \frac{J}{\ell}.
\]

Near this organization-free equilibrium, the process dynamics linearize
to

\[
\dot{x} = (r_0 B - dI)x.
\]

Let \(\rho(B)\) denote the spectral radius of \(B\). The leading growth
rate is

\[
s = r_0 \rho(B) - d,
\]

which motivates the dimensionless organization reproduction number

\[
R_{\rm org} = \frac{J\rho(B)}{\ell d}. \tag{3}
\]

If \(R_{\rm org} < 1\), sufficiently small organization decays. If
\(R_{\rm org} > 1\), it can grow from low abundance. For irreducible
\(B\), let \(p\) be the positive Perron eigenvector normalized by
\(\mathbf{1}^T p = 1\) and let \(\lambda = \rho(B)\). The positive
equilibrium is

\[
r^* = \frac{d}{\lambda}, \qquad x^* = Np, \qquad N = \frac{J - \ell d/\lambda}{d},
\]

which exists precisely when \(R_{\rm org} > 1\).

The same calculation can be written in next-generation-matrix form (van
den Driessche and Watmough, 2002). At the organization-free equilibrium,

\[
F = r_0 B, \qquad V = dI,
\]

so

\[
K = FV^{-1} = \frac{r_0}{d}B, \qquad \rho(K) = \frac{J\rho(B)}{\ell d} = R_{\rm org}.
\]

The vocabulary is borrowed from infectious-disease dynamics, but the
interpretation is generic: \(R_{\rm org}\) asks whether one generation
of executing processes collectively produces more than one replacement
generation before being lost.

For the two-process cycle

\[
B = \begin{pmatrix} 0 & \alpha \\ \beta & 0 \end{pmatrix},
\]

we have \(\rho(B) = \sqrt{\alpha\beta}\). Neither process need reproduce
itself directly. A closed return path across multiple process types can
be enough. If the return path is broken so that the matrix is acyclic,
the spectral radius falls to zero. This isolates the causal object of
interest: not activity alone, but a route from present activity back to
renewed productive capacity.

The positive equilibrium of the base model appears continuously at
\(R_{\rm org} = 1\):

\[
N = \frac{J}{d}\left(1 - \frac{1}{R_{\rm org}}\right).
\]

This is the expected forward transcritical geometry. Establishment and
loss occur at the same threshold in the base model. Section 6 shows how
nonlinear feedback can separate establishment from maintenance; Section
7 then shows a different effect, in which inheritance and continued
reproduction alter successor accessibility even without bistability.

\section{An innovation must pay its opportunity
cost}
A new process is not automatically an improvement. To make this
explicit, begin with \(A \to B\) and \(B \to A\) at unit coefficient.
Let \(A\) divert a fraction \(u\) of its production capacity from \(B\)
to a new process \(C\), while \(C\) contributes back to production of
\(A\) with coefficient \(v\):

\[
B(u,v) = \begin{pmatrix} 0 & 1 & v \\ 1-u & 0 & 0 \\ u & 0 & 0 \end{pmatrix}, \qquad 0 < u < 1,\ v \geq 0.
\]

The dominant eigenvalue is

\[
\lambda(u,v) = \sqrt{1-u+uv} = \sqrt{1+u(v-1)}.
\]

Hence the new process lowers the external supply required for
maintenance only when \(v > 1\). If \(v < 1\), the diversion to \(C\)
reduces collective replacement capacity. The model therefore contains no
rule that rewards novelty, diversity, or complexity. An addition must
causally return enough productive benefit to offset its cost.

This distinction matters for the later argument. The ratchets identified
below are not generated by a preference for larger networks. They are
generated by a change in the dynamical conditions of persistence ---
either because a new feedback loop closes (Sections 5--6), or because
reproduction at scale changes the population of variants available for
the \emph{existing} loop to be locally re-tuned (Section 7).

\section{A seed-dependent extension: history enters the
state}
A fixed linear production matrix does not capture an innovation that
becomes self-maintaining only after it is present. We therefore consider
a nonlinear extension with reactions represented by

\[
R + B \to A + B, \qquad R + A \to A + B, \qquad R + A + C \to A + 2C, \qquad R + C \to A + C.
\]

The corresponding normalized equations are

\[
\dot a = rb + vrc - a, \tag{4}
\] \[
\dot b = ra - b, \tag{5}
\] \[
\dot c = urac - c, \tag{6}
\] \[
\dot r = J - r - (rb + vrc + ra + urac). \tag{7}
\]

All production consumes substrate, and \(\dot T = J - T\) for
\(T = r+a+b+c\).

The \(A\)--\(B\) subsystem can exist without \(C\). Process \(C\),
however, reproduces only when both \(A\) and \(C\) are already present,
while \(C\) contributes back to \(A\). A one-off seed can therefore
activate a loop that was dynamically inaccessible from \(c=0\).

For \(u=30\) and \(v=2\), numerical experiments give the following
controlled comparison. Start the \(A\)--\(B\) system at its equilibrium
with \(J=1.2\). At time \(t=50\), introduce \(10^{-4}\) of \(C\) while
subtracting the same amount from resource. At \(t=150\), lower supply to
\(J=0.9\). With the \(C \to A\) return path intact, \(A\), \(B\), and
\(C\) approach a positive equilibrium. If \(C\) is never seeded, the
organization disappears after the supply reduction. If \(C\) is seeded
but its return effect on \(A\) is removed (\(v=0\)), the organization
also disappears. Thus the persistent difference is carried by the newly
established feedback architecture, not by continuing external supply of
\(C\).

The seed need not be deliberately introduced. One may add a rare local
reaction \(R+A \to A+C\) and treat its occurrence stochastically. Once a
rare event creates \(C\), the ordinary deterministic dynamics can
amplify it. This does not explain why such a reaction must exist in
nature; it shows only that an external evaluator that recognizes and
installs the ``better'' organization is not logically required.

\section{Mechanism I --- retention: the backward bifurcation and
hysteretic
accessibility}
The key result is that the threshold for establishing \(C\) is not the
same as the threshold for maintaining the \(C\)-containing state.

For \(c>0\), the equilibrium conditions of Eq. (7) imply

\[
a = \frac{1}{ur}, \qquad b = \frac{1}{u}, \qquad c = \frac{1-r^2}{uvr^2}, \qquad 0<r<1.
\]

The associated resource supply is

\[
J(r) = r + \frac{1/r + 1 + (1/r^2-1)/v}{u}. \tag{8}
\]

The fold of this positive equilibrium branch satisfies

\[
\frac{dJ}{dr}=0 \iff ur^3 - r - \frac{2}{v} = 0.
\]

For \(u=30\) and \(v=2\),

\[
r_{\rm fold} \approx 0.356236, \qquad J_{\rm fold} \approx 0.597806.
\]

Now examine invasion from the resident \(A\)--\(B\) state. For \(J>1\),
that equilibrium has

\[
r^*=1, \qquad a^*=b^*=\frac{J-1}{2}.
\]

At very small \(c\),

\[
\dot c \approx (ur^*a^*-1)c,
\]

so the type-reproduction number of \(C\) is

\[
R_C = ur^*a^* = u\frac{J-1}{2}.
\]

The invasion threshold \(R_C=1\) is therefore

\[
J_{\rm inv} = 1+\frac{2}{u} = 1.066667. \tag{9}
\]

This is the same logic used for type-specific invasion numbers in
heterogeneous epidemic systems (Heesterbeek and Roberts, 2007).

The striking fact is

\[
J_{\rm fold} < 1 < J_{\rm inv}.
\]

The positive \(C\)-containing equilibrium branch extends backward from
the invasion point to much lower values of \(J\), turning at the fold.
This is the characteristic geometry called a backward bifurcation in
epidemiological models with sub-threshold endemic equilibria
(Kribs-Zaleta and Velasco-Hernández, 2000; van den Driessche and
Watmough, 2002).

The stable states can be summarized as follows:

\begin{longtable}[]{@{}lll@{}}
\toprule\noalign{}
Supply range & Stable unexpanded state & Stable expanded state \\
\midrule\noalign{}
\endhead
\bottomrule\noalign{}
\endlastfoot
\(J < J_{\rm fold}\) & organization-free & no \\
\(J_{\rm fold} < J < 1\) & organization-free & \(A\)--\(B\)--\(C\) \\
\(1 < J < J_{\rm inv}\) & \(A\)--\(B\) & \(A\)--\(B\)--\(C\) \\
\(J > J_{\rm inv}\) & \(A\)--\(B\) unstable to rare \(C\) &
\(A\)--\(B\)--\(C\) \\
\end{longtable}

Thus, at \(J=0.9\), an unexpanded system decays to the organization-free
state while an already established \(A\)--\(B\)--\(C\) organization
persists. At \(J=1.02\), a small \(C\) introduced into the \(A\)--\(B\)
state cannot invade, but an already established \(A\)--\(B\)--\(C\)
state can persist. In both cases the \textbf{external environment is
identical}, yet the long-term outcome differs with realized state and
history. This is hysteresis.

\begin{quote}
An organizational innovation exhibits \textbf{Mechanism I (retention)}
when the conditions required to establish it are stricter than the
conditions required to maintain it once established.
\end{quote}

In terms of the accessibility kernel, Mechanism I does \textbf{not} act
by changing \(E\). Its scientific content is precisely that \(E\) may be
held fixed while \(\mathcal A_T\) differs because the realized state
\(z\) occupies a different basin. The backward bifurcation supplies one
explicit geometry for this state-dependent accessibility. The new loop
changes the persistence landscape; once the expanded state exists,
reversing the environmental change that permitted establishment does not
necessarily remove it. This is stronger than saying that positive
feedback stabilizes a system. It is a claim that realized organization
changes what remains accessible under the same external conditions.

\section{Mechanism II --- inherited adaptive extension and rescue
near
threshold}
Section 6 shows one route by which realized organization changes future
accessibility: the same external environment can support different
persistent outcomes because the realized state occupies a different
basin. Mechanism II is different. It requires neither a backward
bifurcation nor multiple attractors. It arises whenever organization is
reproduced with inheritance and variation.

Two contributions should be distinguished. \textbf{Inheritance changes
the origin of search:} descendants are generated as modifications of an
already viable organization rather than from the organization-free
state. \textbf{Continued reproduction changes the amount of search:}
while the lineage remains extant, it generates a finite number of trials
from which a viable successor may appear. Near a persistence boundary,
these effects combine with the fact that only a small adaptive
displacement may be required to recover positive growth.

\subsection{Establishment changes cumulative exploratory
opportunity}
Consider first a simple subcritical branching process with mean
reproduction number \(R<1\), initiated by one founder. The expected
\textbf{total cumulative progeny}, including the founder, is

\[
E[N_{\rm tot}] = \frac{1}{1-R}. \tag{10}
\]

At \(R=0.5\), \(E[N_{\rm tot}]=2\); at \(R=0.9\), it is \(10\); at
\(R=0.99\), it is \(100\). As \(R\to1^-\), expected cumulative progeny
per founder grows sharply. This does \textbf{not} imply that the
standing population at a given moment is large, nor that standing
variation itself scales as \(1/(1-R)\). It means that a lineage
initiated near criticality is expected to generate more cumulative
reproductive events before extinction than one far below threshold.

This motivates a more general quantity. Let \(N(t)\) denote the
abundance or number of active realizations of an organization, and let
\(\nu(t)\) denote the rate per active realization at which relevant
candidate variants are generated. Define the \textbf{exploratory
opportunity} over horizon \(T\) as

\[
\boxed{
\Omega_T = \int_0^T N(t)\,\nu(t)\,dt.
}
\tag{11}
\]

\(\Omega_T\) is not an independent argument that must be added to
\(\mathcal A_T\); it is an intermediate quantity generated by
\((\theta,z,E,q)\) and helps determine the kernel. In the simple
branching-process case with approximately constant per-reproduction
variant probability, Eq. (10) provides the corresponding scaling of
cumulative opportunity near \(R=1\).

This is the connection emphasized in pathogen-emergence theory: even
while \(R_0<1\), transmission chains become longer on average as the
threshold is approached, creating more opportunities for adaptive
mutations to arise before extinction (Antia et al., 2003; Woolhouse et
al., 2005). The same abstract logic applies to any reproducing
organization for which repeated executions provide opportunities for
inherited modification.

\subsection{Deterioration from above: distance to positive
growth}
Now consider an organization that has already been established and whose
environment deteriorates. Write the resident continuation number as

\[
R_{\rm res}=1-\varepsilon, \qquad 0<\varepsilon<1.
\]

A variant with multiplicative reproductive advantage \(s\) has
approximately

\[
R_{\rm var}=(1-\varepsilon)(1+s).
\]

The variant restores positive growth when

\[
(1-\varepsilon)(1+s)>1
\quad\Longleftrightarrow\quad
s>\frac{\varepsilon}{1-\varepsilon}
\approx \varepsilon
\qquad(\varepsilon\ \text{small}).
\tag{12}
\]

Thus a resident only slightly below its persistence boundary may require
only a small accessible improvement to produce a growing successor,
whereas a resident far below threshold requires a much larger change.
This is a statement about distance to the growth boundary, not about
purpose or effort.

A more general local form makes the geometry explicit. Let
\(R(\theta,E)\) denote continuation of organization \(\theta\) in
environment \(E\), and suppose \(R(\theta,E)=1-\varepsilon\). For a
nearby variant \(\theta'=\theta+\delta\theta\),

\[
R(\theta+\delta\theta,E)
\approx
1-\varepsilon
+
\nabla_\theta R(\theta,E)\cdot\delta\theta.
\]

Whenever locally accessible variation includes a favourable direction,
positive growth is restored if

\[
\nabla_\theta R(\theta,E)\cdot\delta\theta>\varepsilon.
\tag{13}
\]

This identifies a geometric opportunity but not a guarantee. If
accessible variants are constrained, the local gradient is negligible,
adaptation requires crossing a fitness valley, or decline is too rapid,
rescue may fail.

The classical term for successful recovery of a declining lineage
through adaptation before extinction is \textbf{evolutionary rescue}
(Bell, 2017). Its probability depends jointly on initial abundance,
standing variation, variant-generation rates, the distribution of
effects, the size of the adaptive displacement required, and the time
available before extinction. In the present language, the relevant point
is that decline is not instantaneous disappearance: a finite
\(\Omega_T\) remains during which \(\mathcal A_T\) may assign non-zero
probability to a persistent successor.

\subsection{Illustration: an eco-evolutionary feedback
loop}
A biological system that makes this feedback unusually visible is the
emergence and subsequent evolution of a pathogen in an initially naive
host population. Reviews of SARS-CoV-2 evolution describe how increasing
population immunity changed the selective environment experienced by
later lineages, making antigenic novelty and immune escape important
components of relative fitness alongside intrinsic transmission-related
traits (Carabelli et al., 2023). In the present terms, prior circulation
altered the ecological environment; that environment changed the
continuation number of resident lineages; the still-existing viral
population generated further variation; and some descendants occupied
regions of phenotype space with renewed positive growth. This is an
eco-evolutionary feedback because consequences of reproduction alter the
environment that subsequently filters descendants (Ferriere and
Legendre, 2013). It is an illustration of the accessibility logic, not a
claim that the toy reaction model of Sections 2--6 is a mechanistic
model of SARS-CoV-2 evolution.

\subsection{Inheritance changes the origin of
search}
The deeper asymmetry does not depend on rescue. A first viable
organization may require a rare transition through a large configuration
space. Once it exists and is faithfully enough inherited, descendants do
not repeat that search from the organization-free state. They begin from
the realized organization:

\[
\theta \longrightarrow \theta+\delta\theta
\]

rather than

\[
0 \longrightarrow \theta+\delta\theta.
\]

Previously solved viability constraints are therefore part of the
initial condition of subsequent search unless later variation destroys
them. This does not mean every successor is better or more complex. It
means that successor generation is conditioned on what survived.

\[
\boxed{\text{Successful previous organization becomes the starting point for the next search.}}
\]

This is inheritance-mediated path dependence, not hysteresis in the
strict dynamical-systems sense. Its consequence for accessibility is
that the kernel is centred on an already realized organization and is
supplied by a history-dependent amount of exploratory opportunity
\(\Omega_T\).

\subsection{One cumulative ratchet, two mechanisms so far, and one
common
object}
Mechanisms I and II are not two unrelated ``ratchets.'\,' They are two
mathematically distinct ways in which realized organization makes future
accessibility history-dependent.

\begin{longtable}[]{@{}
  >{\raggedright\arraybackslash}p{(\columnwidth - 6\tabcolsep) * \real{0.2500}}
  >{\raggedright\arraybackslash}p{(\columnwidth - 6\tabcolsep) * \real{0.2500}}
  >{\raggedright\arraybackslash}p{(\columnwidth - 6\tabcolsep) * \real{0.2500}}
  >{\raggedright\arraybackslash}p{(\columnwidth - 6\tabcolsep) * \real{0.2500}}@{}}
\toprule\noalign{}
\begin{minipage}[b]{\linewidth}\raggedright
Mechanism
\end{minipage} & \begin{minipage}[b]{\linewidth}\raggedright
What history changes
\end{minipage} & \begin{minipage}[b]{\linewidth}\raggedright
Accessibility effect
\end{minipage} & \begin{minipage}[b]{\linewidth}\raggedright
Formal object
\end{minipage} \\
\midrule\noalign{}
\endhead
\bottomrule\noalign{}
\endlastfoot
I --- retention & basin occupancy and persistence geometry at fixed
\(E\) & what can remain & state \(z\) and basin structure \\
II --- inherited adaptive extension & inherited base point and
cumulative exploratory opportunity & where search begins and how many
trials occur & \(\theta\) and \(\Omega_T\) \\
III --- evolvability & generator of candidate successors & how search is
performed & \(q(\theta'\mid\theta)\) \\
\end{longtable}

Mechanism I can retain an organization under conditions that would not
establish it from rarity. Mechanism II can then allow that retained
organization to generate local successors as conditions continue to
change, without requiring another bistable branch. Their combination
gives

\[
\begin{aligned}
\text{establishment} &\to \text{retention} \to \text{reproduction} \to \text{exploratory opportunity} \\
&\to \text{inherited variation} \to \text{successful successor} \to \text{renewed persistence} \to \cdots
\end{aligned}
\]

Nothing in this sequence implies that the system senses its distance to
a threshold or deliberately increases variation when threatened. The
relevant quantities arise from ordinary reproduction, turnover,
variation, and differential continuation. Near a persistence boundary,
the \textbf{filter} on already generated variation changes because small
positive differences in continuation become consequential. The kernel
may widen, narrow, or collapse; the theory predicts historical
dependence, not guaranteed improvement.

\section{Mechanism III --- evolvability: from ratchets of
organization to a ratchet of
search}
Mechanisms I and II modify future accessibility without changing the
variant-generating mechanism itself: the first through state-dependent
persistence, the second through inherited starting points and
history-dependent exploratory opportunity. Neither establishes that the
system will become progressively better at generating viable successors.
That additional step requires the generator of candidate organizations
itself to vary.

That additional step requires variation in the innovation-generating
architecture itself. Let \(\theta\) denote the current organization ---
for example its production matrix, nonlinear interaction coefficients,
modular structure, connection rules, or learning rules. Let \(m\) denote
a modifier that changes how candidate variants \(\theta'\) are generated
from \(\theta\), represented by a kernel

\[
q_m(\theta' \mid \theta).
\]

A modifier may change mutation rates, recombination, exploratory
connectivity, modularity, the probability that new feedback paths are
formed, the cost of testing variants, or the ability to preserve
successful configurations. The generator \(q_m\) determines the
distribution of candidate successors, but not by itself how many
candidate-generating events occur. The realized history also determines
\(N(t)\) and therefore the exploratory opportunity \(\Omega_T\) of Eq.
(11). Evolvability can thus change accessibility by altering the
distribution of candidate successors, while ecology and persistence
determine how many draws from that distribution occur.

Define a persistence-relevant quantity \(\Phi(\theta)\), such as the
range of external conditions over which the organization remains stable,
the expected duration of persistence under an environment distribution,
or a multi-dimensional measure combining replacement capacity with
resource cost. One possible local measure of innovation yield is

\[
E(m \mid \theta) = \int \left[\Phi(\theta') - \Phi(\theta)\right]_+ q_m(\theta' \mid \theta)\, d\theta' - c(m), \tag{14}
\]

where \(c(m)\) is the direct cost of the modifier. Equation (14) is not
a new biological law; it is a bookkeeping device that separates three
quantities that are otherwise easily conflated: the rate of variation,
the probability that variation improves persistence, and the cost of
generating or testing it.

If a modifier that increases \(E\) is retained together with the
successful organizations it helps generate, then a third feedback
becomes possible:

\[
\begin{aligned}
\text{innovation-generating architecture}
&\to \text{more retainable innovations}\\
&\to \text{greater persistence or reproduction}\\
&\to \text{retention of the architecture.}
\end{aligned}
\]

The system can then, in a precise sense, become better at becoming
differently organized. This is closely related to the established
concept of evolvability: the capacity to generate heritable, selectable
phenotypic variation (Wagner and Altenberg, 1996; Gerhart and Kirschner,
2007). Theory and experiments show that mutation rates, recombination,
robustness, modularity, and genotype--phenotype structure can affect
evolvability. Models have demonstrated conditions under which
evolvability can itself be indirectly selected (Earl and Deem, 2004).

However, this step must not be treated as automatic. Selection for
future adaptive capacity is generally indirect and can be weaker or more
restrictive than direct selection on present performance (Díaz Arenas
and Cooper, 2013; Wagner, 2022). Recent reviews continue to distinguish
clearly between a trait affecting evolvability and selection acting
because of that effect (Pélabon et al., 2025).

This caveat strengthens the present argument. We do not need to assert
that ``systems seek complexity'' or that increased variation is
intrinsically advantageous. More variation can be destructive. Excessive
connectivity can propagate failure. Strong feedback can lock in
maladaptive states. An effective ratchet --- of any of the three kinds
identified above --- can preserve harmful organization as easily as
beneficial organization relative to a chosen external criterion.

Mechanism III is therefore more specific than ``systems get better at
adapting'':

\begin{quote}
Self-maintaining systems can, under conditions that couple
innovation-generating mechanisms to the persistence or reproduction of
the organizations they help create, evolve mechanisms that increase the
rate at which maintainable innovations are generated, tested,
integrated, and retained.
\end{quote}

This gives a hierarchy of causal closure without invoking foresight, in
which Mechanisms I, II, and III occupy successive levels:

\begin{itemize}
\tightlist
\item
  \textbf{Level 1:} processes produce processes.
\item
  \textbf{Level 2:} a network collectively maintains its organization
  (\(R_{\rm org}\)).
\item
  \textbf{Level 3 --- Mechanism I:} new feedback changes the conditions
  of its own persistence, i.e.~the dynamical basin (structural
  hysteresis, Section 6).
\item
  \textbf{Level 4 --- Mechanism II:} reproduction with inheritance
  changes the base point of successor search, while continued activity
  determines the cumulative exploratory opportunity available before
  extinction (inheritance-mediated path dependence coupled to
  evolutionary rescue, Section 7).
\item
  \textbf{Level 5 --- Mechanism III:} the organization changes how new
  variants are generated and retained --- the variation operator itself.
\item
  \textbf{Level 6:} those innovation-generating mechanisms themselves
  vary and are retained differentially.
\end{itemize}

At Level 6, the system can strengthen the process by which it
strengthens itself. Read down this list, each level changes the
finite-horizon accessibility kernel through a different causal channel:
whether organization persists at all; which basin is occupied at fixed
environment; which inherited organization serves as the base point of
variation and how much exploratory opportunity its history supplies; how
candidate successors are generated; and finally how that generator
itself changes. This is a plausible mechanism for cumulative
organization, but the top of the hierarchy remains a hypothesis until
the modifier dynamics are explicitly modelled and tested.

\section{Early-warning signatures of the
thresholds}
The framework also makes dynamical predictions about loss of resilience.
Near a stable equilibrium \(z^*\), add weak stochastic forcing and
linearize:

\[
\dot{\delta z} = A(J)\delta z + \sigma \xi(t).
\]

Projecting onto the slow eigenmode with eigenvalue
\(\lambda_{\rm slow}(J)<0\) yields an Ornstein--Uhlenbeck approximation.
Under standard assumptions,

\[
{\rm Var}(\delta z) \propto \frac{1}{|\lambda_{\rm slow}|}, \qquad \rho(\Delta t) \approx 1 - |\lambda_{\rm slow}|\Delta t.
\]

Thus, as the dominant restoring eigenvalue approaches zero, recovery
slows, variance increases, and lagged autocorrelation rises: the
standard critical-slowing-down picture (Scheffer et al., 2009; Dakos et
al., 2012).

The two thresholds of Section 6 have different local scalings. For the
base transcritical threshold,

\[
|\lambda_{\rm slow}| \propto |R_{\rm org}-1|, \qquad {\rm Var} \propto |R_{\rm org}-1|^{-1}.
\]

For the stable \(C\)-containing branch approaching the fold from above,

\[
|\lambda_{\rm slow}| \propto \sqrt{J - J_{\rm fold}}.
\]

Numerical evaluation of the stable branch gives a fitted exponent of
approximately 0.515, consistent with the generic square-root scaling of
a saddle-node. Hence

\[
{\rm Var} \propto (J-J_{\rm fold})^{-1/2}.
\]

Two corrections are worth making explicit. First, the stable branch must
be used for the return-to-equilibrium interpretation; the opposite side
of the fold is unstable. Second, \(1/\sqrt{\Delta J}\) diverges less
strongly than \(1/\Delta J\) as \(\Delta J \to 0\). Therefore the fold
does not imply a stronger variance divergence at the same numerical
distance to threshold. What differs is the local scaling law and, more
generally, the geometry of approach to the transition.

Mechanism II adds a third, statistical rather than spectral, signature.
Equation (10) predicts expected cumulative progeny in a simple
subcritical branching process, not instantaneous standing population.
More generally, Eq. (11) motivates measuring exploratory opportunity
\(\Omega_T\): the cumulative number of relevant variant-generating
trials available while the resident remains extant. All else equal,
lineages nearer the persistence boundary can combine a larger expected
cumulative reproductive opportunity per founder with a smaller adaptive
displacement required for recovery. Empirical tests should therefore
manipulate or estimate both threshold distance and \(\Omega_T\), rather
than treating proximity to \(R=1\) as a proxy for population size.

These predictions are falsifiable but should not be oversold. Empirical
early-warning signals are sensitive to noise structure, sampling,
nonstationarity, rates of parameter change, and the observability of the
slow mode. The equations predict what this model would do; they do not
establish that a real biological, social, technological, or AI system is
governed by the same local normal form.

\section{What the mechanism explains --- and what it does
not}
The construction supports six claims.

\begin{enumerate}
\def\labelenumi{\arabic{enumi}.}
\tightlist
\item
  \textbf{Collective self-maintenance is possible without a central
  controller.} Local production processes can jointly replace the
  processes that execute the network.
\item
  \textbf{Self-maintenance has a threshold.} The base model has a
  reproduction number \(R_{\rm org}\) that separates decay from
  establishment at low abundance.
\item
  \textbf{An innovation can alter its own persistence conditions.} A new
  feedback loop can produce bistability and hysteresis, making the
  innovation easier to maintain than to establish (structural
  hysteresis).
\item
  \textbf{An established organization can have greater access to viable
  successors than a de novo system, even without any new bistable
  feedback loop.} Inheritance supplies an already-functional starting
  point; continued reproduction supplies a finite exploratory
  opportunity \(\Omega_T\); and proximity to the persistence boundary
  can reduce the adaptive displacement needed to restore positive
  growth.
\item
  \textbf{A temporary event can become embodied in persistent
  organization.} The state of the network can therefore carry a minimal
  form of history or memory.
\item
  \textbf{There is a logically coherent route to third-order
  improvement.} If mechanisms that generate and retain variants
  themselves vary and remain coupled to the persistence of their
  descendants, evolvability can become part of the feedback structure.
\end{enumerate}

The construction does not show that:

\begin{itemize}
\tightlist
\item
  complexity must increase monotonically;
\item
  every added process is useful;
\item
  every persistent process is useful for the whole network;
\item
  more variation, connectivity, or feedback is always advantageous;
\item
  backward bifurcation is necessary for all evolutionary ratchets, or
  that it is the only mechanism the term should be reserved for;
\item
  inherited adaptive extension or evolutionary rescue will succeed
  whenever an organization approaches its threshold --- the opportunity
  window can close before a sufficient successor is generated (Bell,
  2017);
\item
  evolvability will always be directly or indirectly selected;
\item
  the present toy reactions are realized in any particular organism,
  ecosystem, institution, robot, or AI system;
\item
  the number of process types is a sufficient measure of biological or
  informational complexity.
\end{itemize}

Indeed, the same model class contains simplification: a sufficiently
fast self-reproducing competitor can displace a multi-process network,
and a resident pushed far below threshold before deterioration is caught
will decay regardless of how large its standing population once was. The
framework therefore supplies mechanisms for cumulative organization, not
a theorem that cumulative organization is inevitable.

\section{A research programme for cumulative
organization}
The argument becomes scientifically useful when translated into
measurements. For a candidate real system, one should identify:

\begin{enumerate}
\def\labelenumi{\alph{enumi})}
\tightlist
\item
  the processes or process classes whose continuation is at stake;
\item
  the resources consumed in their replacement;
\item
  the directed causal effects by which one process increases the future
  abundance or execution of another;
\item
  the replacement or invasion threshold for the existing organization;
\item
  a candidate innovation and its opportunity cost;
\item
  whether the innovation creates a return path into the organization
  that produces or maintains it;
\item
  whether the threshold for invasion differs from the threshold for loss
  once established (structural hysteresis);
\item
  whether current state depends on historical seeding or initial
  conditions under identical external conditions;
\item
  how cumulative exploratory opportunity \(\Omega_T\) changes with
  abundance, turnover, variant-generation rate, and distance from the
  persistence boundary;
\item
  whether recovery from deterioration depends jointly on \(\Omega_T\),
  inherited starting structure, and the adaptive displacement required
  to regain positive growth;
\item
  whether variation-generating or connection-generating mechanisms
  themselves differ between persistent lineages or organizations;
\item
  whether those mechanisms are retained because of the innovations they
  make more likely, rather than because of an unrelated direct benefit.
\end{enumerate}

The decisive experiments separate the mechanisms. For retention,
introduce the same candidate process in otherwise matched systems,
preserve resource cost, and selectively break the proposed return path;
if the persistence advantage remains, the loop was not causal. For
inherited adaptive extension, manipulate independently (i) the inherited
starting organization, (ii) distance from the persistence boundary, and
(iii) cumulative exploratory opportunity \(\Omega_T\). This
distinguishes a true history-dependent search advantage from a simple
effect of current abundance. Likewise, for the third-order hypothesis,
compare architectures that differ in their innovation operator while
controlling direct present-day benefits. The central empirical object is
not ``complexity'' in the abstract but the measurable causal
architecture by which present consequences alter future production and
future variation.

\section{Discussion: from Darwinian selection to the persistence of
organization}
Darwinian natural selection explains differential propagation of
heritable variation. The present framework addresses a complementary
dynamical question, stated already in Section 1.1 and worth restating
here as the paper's real subject: not how the present state influences
the next state, which any dynamical system does, but \textbf{how
realized organization can endogenously change the finite-horizon
transition kernel over viable successor organizations}. Mechanisms I
through III are three concrete, mechanistic answers rather than three
unrelated case studies: state-dependent persistence changes what can
remain; inheritance and continued reproduction change where successor
search begins and how many trials occur; and evolvability changes the
generator of candidate successors itself.

The answer offered here does not require a force pulling matter toward
complexity. Open systems receive resource and undergo local
transformations. Some configurations fail to replace themselves. Others
close productive return paths and persist. Some new processes impose
costs and disappear. Others alter the network such that the expanded
state becomes self-supporting across a wider region of parameter space
--- Mechanism I. Independently of whether any such new loop has closed,
an organization that has been reproducing above threshold for some time
carries with it a population of reproducing realizations and any
standing variation; if deterioration brings its own reproduction back
toward --- rather than past --- its threshold, only a small correction,
evaluated against its already-functional configuration rather than
against the empty state, is needed to recover growth --- Mechanism II.
Both mechanisms make history dynamically consequential, but through
different causal channels in \(\mathcal A_T\): one changes persistence
through state-dependent basin structure at fixed environment; the other
changes the inherited origin of variation and the cumulative opportunity
available to generate successors.

This also clarifies why ``more feedback'' is not itself the objective,
and why ``more reproduction'' is not either. What matters for Mechanism
I is feedback that survives its costs and changes the continuation of
the organization that generates it. What matters for Mechanism II is the
conjunction of an inherited viable base point, a non-zero exploratory
opportunity, and a selection filter under which small improvements can
become decisive near a persistence boundary --- not a preference or
incentive of the system itself. No system need represent a future goal
for either mechanism to operate, and neither guarantees that
\(\mathcal{A}\) widens: accessibility can just as easily narrow,
trapping a lineage on whatever it has already found, as it can expand.

If, in addition, architectures differ in the distribution of variants
they generate, the same logic can recurse a further level: an
architecture can be retained partly because it generates descendants or
successor states that more often acquire maintainable improvements, or
that more often survive threshold crossings via rescue --- Mechanism
III, acting on \(q\) itself. This recursive step is where the argument
intersects the literature on evolvability. Wagner and Altenberg
emphasized that evolution depends not only on variation but on the
structure mapping variation into phenotypic consequences, and
highlighted modularity as one route by which evolvability may change
(Wagner and Altenberg, 1996). Gerhart and Kirschner's
facilitated-variation framework similarly emphasizes conserved core
processes whose organization can make useful phenotypic change more
accessible (Gerhart and Kirschner, 2007) --- a claim closely related to
the inheritance argument of Section 7.4, in which conserved organization
is precisely what makes the \emph{next} useful change a local rather
than a global search. Earl and Deem demonstrated in explicit models that
evolvability can be selected under changing environments (Earl and Deem,
2004). At the same time, experimental and theoretical reviews warn that
indirect selection for evolvability is conditional and can be confused
with direct selection on traits that merely have evolvability as a side
effect (Díaz Arenas and Cooper, 2013; Wagner, 2022; Pélabon et al.,
2025).

The proposed hierarchy therefore does not replace Darwinian selection.
It extends the object to which selection and dynamical persistence can
apply: not only traits, but the network architecture that generates
traits; not only an organization, but its capacity to recover a viable
form of itself, and beyond that, its capacity to produce and retain new
organization altogether.

A concise way to state the combined result is:

\begin{quote}
A self-maintaining network can retain a new organization when that
organization contributes to the processes that reproduce it. Once
established, inheritance makes that organization the base point for
successor search, while continued reproduction supplies a finite
exploratory opportunity from which locally viable successors may arise.
If the retained organization also changes how future variants are
generated, connected, tested, or preserved, then the search process
itself can enter the same feedback loop.
\end{quote}

Underneath the three mechanisms is one abstract operation:

\[
\boxed{\text{Realized organization transforms the accessibility of organization that can follow it.}}
\]

This is deliberately narrower than saying that the present state
influences the next state, a claim true of essentially every dynamical
system. The scientific content is that realized organization changes the
\textbf{transition structure over viable successors}. Mechanism I does
so through basin occupancy and persistence geometry at fixed \(E\);
Mechanism II through the inherited base point and the history-dependent
exploratory opportunity \(\Omega_T\); Mechanism III through the
generator \(q\) of candidate successors. None requires foresight.

Because the operation is descriptive rather than normative, ``the
ratchet'' in this paper does not mean a mechanism that only ever helps.
It means: \textbf{a cumulative organizational ratchet exists when
realized organization is retained in a way that changes the
accessibility distribution of subsequent organization, so that
previously achieved viability need not be rediscovered de novo} (Section
1.1). That definition is silent on direction. Cumulative organization
--- the case this paper is ultimately interested in --- is the special
case in which the transformation happens to preserve enough previously
viable organization that later search need not repeatedly start from
zero. It is not the only possible outcome of the same operation.

This is sufficient to make cumulative emergence possible. It is not
sufficient to make it inevitable.

\section{Conclusion}
We have combined four layers of one argument around a single object: the
finite-horizon accessibility kernel \(\mathcal A_T(S\mid\theta,z,E,q)\)
over persistent successor organizations.

First, an open network of local production processes can collectively
cause its own continuation. Its establishment is governed by a
reproduction threshold, \(R_{\rm org}\). This determines whether
organization is realized at all.

Second (\textbf{Mechanism I --- retention}), a seed-dependent innovation
can create a feedback loop that separates establishment from
maintenance. In the worked example, the innovation invades only above
\(J_{\rm inv}=1.0667\) but persists once established down to
\(J_{\rm fold}=0.5978\). The resulting backward bifurcation produces
bistability and hysteresis. At identical external conditions, different
realized states occupy different basins and therefore have different
persistence accessibility. Realized organization changes \textbf{what
can remain}.

Third (\textbf{Mechanism II --- inherited adaptive extension}), no
bistability is required. Reproduction with inheritance makes an already
viable organization the base point of subsequent variation, while
continued activity generates a cumulative exploratory opportunity
\(\Omega_T\). For a simple subcritical branching process, expected total
progeny per founder scales as \(1/(1-R)\) as \(R\to1^-\); this is a
statement about cumulative reproductive opportunity, not instantaneous
standing population. When a resident is only \(\varepsilon\) below its
persistence boundary, a multiplicative advantage
\(s>\varepsilon/(1-\varepsilon)\approx\varepsilon\) is sufficient to
restore positive growth. Realized organization therefore changes
\textbf{where search begins, how many trials can occur, and how far a
locally accessible successor may need to move to recover viability}.

Fourth (\textbf{Mechanism III --- evolvability}), the same logic can
recurse to the generator of variation itself. If architectures differ in
how they generate, connect, test, or preserve candidate successors, and
those differences remain coupled to descendant persistence, selection
can alter \(q\). Realized organization then changes \textbf{how the next
search is performed}.

The central claim is neither that nature seeks complexity nor that
complexity must increase. Accessibility may widen, narrow, or collapse;
inherited organization may preserve liabilities as readily as useful
structure. The claim is more basic: realized organization can change the
transition kernel over what can follow it. Cumulative organization is
the special case in which those history-dependent transformations
preserve enough earlier viability that subsequent viable organization
need not repeatedly be reconstructed from scratch.

\[
\boxed{\text{Existence changes what existence can follow.}}
\]

Under conditions of retained viability, this yields the more specific
cumulative principle:

\[
\boxed{\text{What persists becomes part of the starting conditions for what can persist next.}}
\]

\section{References}
Antia, R., Regoes, R. R., Koella, J. C. and Bergstrom, C. T. (2003). The
role of evolution in the emergence of infectious diseases.
\emph{Nature}, 426, 658--661. \\url{https://doi.org/10.1038/nature02104}.

Bell, G. (2017). Evolutionary rescue. \emph{Annual Review of Ecology,
Evolution, and Systematics}, 48, 605--627.
\\url{https://doi.org/10.1146/annurev-ecolsys-110316-023011}.

Carabelli, A. M., Peacock, T. P., Thorne, L. G., Harvey, W. T., Hughes,
J., COVID-19 Genomics UK Consortium et al.~(2023). SARS-CoV-2 variant
biology: immune escape, transmission and fitness. \emph{Nature Reviews
Microbiology}, 21, 162--177. \\url{https://doi.org/10.1038/s41579-022-00841-7}.

Dakos, V., Carpenter, S. R., Brock, W. A., Ellison, A. M., Guttal, V.,
Ives, A. R., Kéfi, S., Livina, V., Seekell, D. A., van Nes, E. H. and
Scheffer, M. (2012). Methods for detecting early warnings of critical
transitions in time series illustrated using simulated ecological data.
\emph{PLoS ONE}, 7(7), e41010. \\url{https://doi.org/10.1371/journal.pone.0041010}.

Díaz Arenas, C. and Cooper, T. F. (2013). Mechanisms and selection of
evolvability: experimental evidence. \emph{FEMS Microbiology Reviews},
37(4), 572--582. \\url{https://doi.org/10.1111/1574-6976.12008}.

Earl, D. J. and Deem, M. W. (2004). Evolvability is a selectable trait.
\emph{Proceedings of the National Academy of Sciences}, 101(32),
11531--11536. \\url{https://doi.org/10.1073/pnas.0404656101}.

Ferriere, R. and Legendre, S. (2013). Eco-evolutionary feedbacks,
adaptive dynamics and evolutionary rescue theory. \emph{Philosophical
Transactions of the Royal Society B}, 368(1610), 20120081.
\\url{https://doi.org/10.1098/rstb.2012.0081}.

Gerhart, J. and Kirschner, M. (2007). The theory of facilitated
variation. \emph{Proceedings of the National Academy of Sciences},
104(Suppl. 1), 8582--8589. \\url{https://doi.org/10.1073/pnas.0701035104}.

Heesterbeek, J. A. P. and Roberts, M. G. (2007). The type-reproduction
number \(T\) in models for infectious disease control.
\emph{Mathematical Biosciences}, 206(1), 3--10.
\\url{https://doi.org/10.1016/j.mbs.2004.10.013}.

Hordijk, W., Hein, J. and Steel, M. (2010). Autocatalytic sets and the
origin of life. \emph{Entropy}, 12(7), 1733--1742.
\\url{https://doi.org/10.3390/e12071733}.

Hordijk, W., Steel, M. and Kauffman, S. (2012). The structure of
autocatalytic sets: evolvability, enablement, and emergence. \emph{Acta
Biotheoretica}, 60(4), 379--392. \\url{https://doi.org/10.1007/s10441-012-9165-1}.

Jain, S. and Krishna, S. (1998). Autocatalytic sets and the growth of
complexity in an evolutionary model. \emph{Physical Review Letters}, 81,
5684--5687. \\url{https://doi.org/10.1103/PhysRevLett.81.5684}.

Jain, S. and Krishna, S. (2001). A model for the emergence of
cooperation, interdependence, and structure in evolving networks.
\emph{Proceedings of the National Academy of Sciences}, 98(2), 543--547.
\\url{https://doi.org/10.1073/pnas.98.2.543}.

Kolchinsky, A. (2025). Thermodynamics of Darwinian selection in
molecular replicators. \emph{Philosophical Transactions of the Royal
Society B}, 380(1936), 20240436. \\url{https://doi.org/10.1098/rstb.2024.0436}.

Kribs-Zaleta, C. M. and Velasco-Hernández, J. X. (2000). A simple
vaccination model with multiple endemic states. \emph{Mathematical
Biosciences}, 164(2), 183--201. \\url{https://doi.org/10.1016/S0025-5564(00)00003-1}.

Pélabon, C. et al.~(2025). Evolvability: progress and key questions.
\emph{BioScience}, 75(12), 1042--1057. \\url{https://doi.org/10.1093/biosci/biaf111}.

Peng, Z., Linderoth, J. and Baum, D. A. (2022). The hierarchical
organization of autocatalytic reaction networks and its relevance to the
origin of life. \emph{PLoS Computational Biology}, 18(9), e1010498.
\\url{https://doi.org/10.1371/journal.pcbi.1010498}.

Scheffer, M., Bascompte, J., Brock, W. A., Brovkin, V., Carpenter, S.
R., Dakos, V., Held, H., van Nes, E. H., Rietkerk, M. and Sugihara, G.
(2009). Early-warning signals for critical transitions. \emph{Nature},
461, 53--59. \\url{https://doi.org/10.1038/nature08227}.

van den Driessche, P. and Watmough, J. (2002). Reproduction numbers and
sub-threshold endemic equilibria for compartmental models of disease
transmission. \emph{Mathematical Biosciences}, 180, 29--48.
\\url{https://doi.org/10.1016/S0025-5564(02)00108-6}.

Wagner, A. (2022). Adaptive evolvability through direct selection
instead of indirect, second-order selection. \emph{Journal of
Experimental Zoology Part B: Molecular and Developmental Evolution},
338, 395--404. \\url{https://doi.org/10.1002/jez.b.23071}.

Wagner, G. P. and Altenberg, L. (1996). Complex adaptations and the
evolution of evolvability. \emph{Evolution}, 50(3), 967--976.
\\url{https://doi.org/10.1111/j.1558-5646.1996.tb02339.x}.

Woolhouse, M. E. J., Haydon, D. T. and Antia, R. (2005). Emerging
pathogens: the epidemiology and evolution of species jumps. \emph{Trends
in Ecology \& Evolution}, 20(5), 238--244.
\\url{https://doi.org/10.1016/j.tree.2005.02.009}.

\end{document}
